\chapter{1996 Nobel Prize in Economics}
\textbf{Laureates: }James Mirrlees (UK), William Vickrey (Canada/USA)

\section*{Reason for Award}
Under asymmetric information incentive economic theory.
\section{What They Did}
James Mirrlees solved optimal taxation problem addressing hidden actions moral hazard problem demonstrating how taxation must balance revenue needs equity incentives distortion Vickrey pioneered auction theory analyzing bidding strategies in sealed-bid auctions revealing optimal reserve prices revenue equivalence principle Vickrey studied mechanism design under asymmetric information where principal cannot observe agent actions or types Vickrey introduced truthful revelation mechanism encouraging participants reveal true valuations His second-price auction design ensures bidding equal to true value dominates strategies Vickrey analysis laid groundwork modern auction theory mechanism design Vickrey-Mirrlees approach focused incentive compatibility ensuring participants truthfully reveal private information necessary for efficient allocation Designing contracts regulatory frameworks accommodating informational asymmetries emerged from their complementary works Vickrey Mirrlees established fundamental principles governing economic transactions under information problems providing unified framework designing institutions mitigating adverse selection moral hazard problems

\section{Impact on Economic Thinking}
Mirrlees Vickrey fundamentally transformed understanding incentive problems arising asymmetric information Optimal taxation theory demonstrates tradeoff between redistribution efficiency losses distortion caused taxation Higher top marginal rates reduce inequality but discourage effort provision Vickrey s auction analysis revealing efficient allocation mechanisms private valuations Second-price auction dominant strategy bidding truthful valuations Revenue equivalence connecting different auction formats informing auction design Vickrey-Mirrlees framework informed regulatory contracting insurance design executive compensation procurement auctions spectrum licensing Mechanism design emerged major field building on their insights Principal-agent problems addressed agency contracts align interests parties Vickrey s truthful revelation inspiring Myerson-Satterthwaite theorem impossibility efficient trade bilateral trading Vickrey Mirrlees established informational asymmetries require explicit attention designing institutions Their work provided systematic framework analyzing incentive compatibility participation constraints efficiency-loss tradeoffs fundamental constraint designing economic organizations institutions facing information limitations Modern contract theory regulatory economics public mechanism design all rest foundations laid Mirrlees Vickrey pioneering insights.
\section{Modern Perspectives Today}
Mechanism design continues expanding Vickrey-Mirrlees foundation Auction formats adapt combinatorial complexities multiple items spectrum auctions energy markets electricity wholesale markets kidney matching school assignment algorithms Vickrey-inspired truthful revelation mechanisms deployed various allocation settings Optimal taxation literature extends Mirrlees framework incorporating behavioral considerations administrative costs Progressive tax design debates mirror Mirrlees efficiency-equity tradeoffs Executive compensation contracts use agency theory reflecting Vickrey-Mirrlees insights Insurance designs accommodate adverse selection Procurement mechanisms incentivize truthful reporting Market design applications Vickrey-Mirrlees insights proliferate Vickrey s truthful revelation inspired mechanism designers create systems eliciting private information efficiently Mirrlees optimal taxation framework informs progressive tax policy discussions Vickrey-Mirrlees principles guide designing institutions coping with asymmetric information their insights remain vital modern institutional design Vickrey-Mirrlees framework continues inspire research exploring mechanisms achieving efficiency under information constraints
